\section{Discussion and conclusions}\label{sec:conclusions}

Joint support cuts are the strongest linear inequalities for a block of
shared nonlinear expressions, and they can be generated with numerical
tools and still be certified. The certificate does not need to cover the
search that found a direction; it covers a lower bound for the exported
coefficients on the whole domain, the exact elimination of the model rows
and the row that the solver stores. Stated in the original variables, the
cuts express the projection of the lifted joint-graph relaxation exactly,
without new variables, and the same analysis shows where this relaxation
stops: rows that enter only through averages leave a Lagrangian duality
gap that no choice of directions removes.

For quadratic blocks the support problem can be solved exactly in two
useful cases, small polytopes and stars with center--leaf rows, and the
second case is solved in nearly linear time. The analysis of paths
explains why such blocks can be stronger than their pairs: pair hulls
glued on finitely many moments of a shared variable accept distributions
that no joint distribution reproduces, exactly when the local zero sets
alternate often enough. The same analysis delimits the gain. A dense
moment relaxation with the products that do not occur in the model closes
the gap on the whole path family, so joint support cuts are best viewed
as a way to obtain that strength in the coordinates the model already has.

The computations temper these conclusions. On MINLPLib models the cuts
were valid but did not change which models SCIP solved, and our Python
implementation made the solves slower. Two reasons stand out. First, SCIP
already obtains much of the joint information by other means; on the path
example its presolve turns the leaves into binaries and strong branching
closes the gap at the root. Second, where the structure occurs, the
frozen work limits allow few cuts, and discovery in Python costs more than
the cuts return on small models. On the constructed family, where SCIP's
root relaxation is provably weak, the cuts close part of the root gap and
reduce the number of nodes. We see three directions in which the
components could become useful in practice: as the certified cut source
of a solver that must justify every relaxation step, extending exact
solving beyond MILP; inside a native constraint handler with discovery
done once, at negligible cost; and for model classes, such as pooling and
networks with shared flows, whose blocks are stars with linking rows.

Several questions remain open. Exact support for trees deeper than stars
with edge-local rows would extend the star algorithm to the structures of
network models; Remark~\ref{rem:trees} shows that the rational partition
argument does not extend directly. A direction search with the guarantees
of Section~\ref{sec:separation} and the speed of a solver callback is not
available, and the complete procedures are far too slow for routine use.
Finally, an activation rule that predicts when joint cuts pay off would be
needed before any default use; our structural rule did not achieve this.
